%%% Hlavní soubor. Zde se definují základní parametry a odkazuje se na ostatní části. %%%

%% Verze pro jednostranný tisk:
% Okraje: levý 40mm, pravý 25mm, horní a dolní 25mm
% (ale pozor, LaTeX si sám přidává 1in)
\documentclass[12pt,a4paper]{report}
\setlength\textwidth{145mm}
\setlength\textheight{247mm}
\setlength\oddsidemargin{15mm}
\setlength\evensidemargin{15mm}
\setlength\topmargin{0mm}
\setlength\headsep{0mm}
\setlength\headheight{0mm}
% \openright zařídí, aby následující text začínal na pravé straně knihy
\let\openright=\clearpage

%% Pokud tiskneme oboustranně:
% \documentclass[12pt,a4paper,twoside,openright]{report}
% \setlength\textwidth{145mm}
% \setlength\textheight{247mm}
% \setlength\oddsidemargin{15mm}
% \setlength\evensidemargin{0mm}
% \setlength\topmargin{0mm}
% \setlength\headsep{0mm}
% \setlength\headheight{0mm}
% \let\openright=\cleardoublepage

%% Použité kódování znaků: obvykle latin2, cp1250 nebo utf8:
%\usepackage[utf8]{inputenc}
%\usepackage[T2A]{fontenc}
%\usepackage[russian,czech,english]{babel}

%\usepackage{fontspec}

%\usepackage{xunicode,fontspec,xltxtra}
%\usepackage[english]{polyglossia}
%\setotherlanguages{russian} % set as "other" so English hyphenation active


\usepackage{libertine}
\setmonofont{Anonymous Pro}

%% Ostatní balíčky
\usepackage{verbatim}
\usepackage{graphicx}
%\usepackage{amsthm}

%% Balíček hyperref, kterým jdou vyrábět klikací odkazy v PDF,
%% ale hlavně ho používáme k uložení metadat do PDF (včetně obsahu).
%% POZOR, nezapomeňte vyplnit jméno práce a autora.
\usepackage[unicode]{hyperref}   % Musí být za všemi ostatními balíčky
\hypersetup{pdftitle=A Comparison of Methods of Czech-to-Russian Machine Translation}
\hypersetup{pdfauthor=Karel Bílek}

\usepackage{tabularx}

\usepackage{polyglossia}
\setmainlanguage{english}
\setotherlanguage{czech}

\usepackage{float}

\usepackage[
   backend=bibtex8      % if we want unicode 
  ,style=iso-authoryear % or iso-numeric for numeric citation method          
  ,babel=other        % to support multiple languages in bibliography
  ,sortlocale=cs_CZ   % locale of main language, it is for sorting
  %,bibencoding=UTF8   % this is necessary only if bibliography file is in different encoding than main document
]{biblatex}


\bibliography{library}
%%% Drobné úpravy stylu
% Tato makra přesvědčují mírně ošklivým trikem LaTeX, aby hlavičky kapitol
% sázel příčetněji a nevynechával nad nimi spoustu místa. Směle ignorujte.
\makeatletter
\def\@makechapterhead#1{
  {\parindent \z@ \raggedright \normalfont
   \Huge\bfseries \thechapter. #1
   \par\nobreak
   \vskip 20\p@
}}
\def\@makeschapterhead#1{
  {\parindent \z@ \raggedright \normalfont
   \Huge\bfseries #1
   \par\nobreak
   \vskip 20\p@
}}
\makeatother

% Toto makro definuje kapitolu, která není očíslovaná, ale je uvedena v obsahu.
\def\chapwithtoc#1{
\chapter*{#1}
\addcontentsline{toc}{chapter}{#1}
}


\newfloat{tabulka}{tbph}{lop}
\floatname{tabulka}{Table}

\long\def\centergraf#1 {
{
    \begingroup
    \centering
    \par
    #1
    \par
    \endgroup
}

}

\def\uv#1{``#1''}

\long\def\jednatabulkan#1#2#3#4{
 \begin{tabulka}[H]
     
         \centergraf{
         \footnotesize
    \begin{tabular}{ #2 }
        #3
        
    \end{tabular}
             
             }
     
     \caption{#4} 
     \label{tabulka:#1}
 \end{tabulka}
}


\long\def\jednatabulkam#1#2#3#4{
 \begin{tabulka}[h]
     
         \centergraf{
         \scriptsize
    \begin{tabularx}{\textwidth}{ #2 }
        #3
        
    \end{tabularx}
             
             }
     
     \caption{#4} 
     \label{tabulka:#1}
 \end{tabulka}
}

\long\def\jednatabulka#1#2#3#4{
 \begin{tabulka}[h]
     
         \centergraf{
         \footnotesize
    \begin{tabularx}{\textwidth}{ #2 }
        #3
        
    \end{tabularx}
             
             }
     
     \caption{#4} 
     \label{tabulka:#1}
 \end{tabulka}
}


\newenvironment{pitemize}{
\begin{itemize}
  \setlength{\itemsep}{5pt}
  \setlength{\parskip}{0pt}
  \setlength{\parsep}{0pt}
}{\end{itemize}}


\long\def\priklady#1 {
    \begin{pitemize}
#1

    \end{pitemize}
}

\long\def\priklad#1#2#3 {
 \item \textbf{#1} \par\nobreak #2 \par\nobreak \emph{#3}

}

\begin{document}

% Trochu volnější nastavení dělení slov, než je default.
\lefthyphenmin=2
\righthyphenmin=2

%%% Titulní strana práce

\pagestyle{empty}
\begin{center}

\large

Charles University in Prague

\medskip

Faculty of Mathematics and Physics

\vfill

{\bf\Large MASTER THESIS}

\vfill

\centerline{\mbox{\includegraphics[width=60mm]{../img/logo.pdf}}}

\vfill
\vspace{5mm}

{\LARGE Karel Bílek}

\vspace{15mm}

% Název práce přesně podle zadání
{\LARGE\bfseries A Comparison of Methods of Czech-to-Russian Machine Translation}

\vfill

% Název katedry nebo ústavu, kde byla práce oficiálně zadána
% (dle Organizační struktury MFF UK)
%Ústav formální a aplikované lingvistiky
Institute of Formal and Applied Linguistics

\vfill

\begin{tabular}{rl}

Supervisor of the master thesis: & doc. RNDr. Vladislav Kuboň, Ph.D. \\
\noalign{\vspace{2mm}}
Study programme: & Informatika  ????  \\
\noalign{\vspace{2mm}}
Specialization: & Asi něco s lingvistikou \\
\end{tabular}

\vfill

% Zde doplňte rok
Prague 2014

\end{center}

\newpage

%%% Následuje vevázaný list -- kopie podepsaného "Zadání diplomové práce".
%%% Toto zadání NENÍ součástí elektronické verze práce, nescanovat.

%%% Na tomto místě mohou být napsána případná poděkování (vedoucímu práce,
%%% konzultantovi, tomu, kdo zapůjčil software, literaturu apod.)

\openright

\noindent
Dedication. I love everybody. (TODO)

\newpage

%%% Strana s čestným prohlášením k diplomové práci

\vglue 0pt plus 1fill

\noindent
I declare that I carried out this master thesis independently, and only with the cited
sources, literature and other professional sources.

\medskip\noindent
I understand that my work relates to the rights and obligations under the Act No.
121/2000 Coll., the Copyright Act, as amended, in particular the fact that the Charles
University in Prague has the right to conclude a license agreement on the use of this
work as a school work pursuant to Section 60 paragraph 1 of the Copyright Act.

\vspace{10mm}

\hbox{\hbox to 0.5\hsize{%
In ........ date ............
\hss}\hbox to 0.5\hsize{%
signature of the author
\hss}}

\vspace{20mm}
\newpage

%%% Povinná informační strana diplomové práce

\vbox to 0.5\vsize{
\setlength\parindent{0mm}
\setlength\parskip{5mm}

Název práce:
Porovnáni metod česko-ruského automatického překladu

% přesně dle zadání

Autor:
Karel Bílek

Katedra:  % Případně Ústav:
Ústav formální a aplikované lingvistiky
% dle Organizační struktury MFF UK

Vedoucí diplomové práce:
doc. RNDr. Vladislav Kuboň, Ph.D.,
Ústav formální a aplikované lingvistiky
%Jméno a příjmení s tituly, pracoviště
% dle Organizační struktury MFF UK, případně plný název pracoviště mimo MFF UK

Abstrakt:
Abstrakt abstrakt.
% abstrakt v rozsahu 80-200 slov; nejedná se však o opis zadání diplomové práce

Klíčová slova:
% 3 až 5 klíčových slov
Čeština, ruština, strojový překlad

\vss}\nobreak\vbox to 0.49\vsize{
\setlength\parindent{0mm}
\setlength\parskip{5mm}

Title:
% přesný překlad názvu práce v angličtině
{A Comparison of Methods of Czech-to-Russian Machine Translation}

Author:
Karel Bílek
%Jméno a příjmení autora

Department:
Institute of Formal and Applied Linguistics
%Ústav formální a aplikované lingvistiky
%Institute of Formal and Applied Linguistics
%Název katedry či ústavu, kde byla práce oficiálně zadána
% dle Organizační struktury MFF UK v angličtině

Supervisor:
doc. RNDr. Vladislav Kuboň, Ph.D.,
Institute of Formal and Applied Linguistics
%Jméno a příjmení s tituly, pracoviště
% dle Organizační struktury MFF UK, případně plný název pracoviště
% mimo MFF UK v angličtině

Abstract:
Abstrakt abstrakt.
% abstrakt v rozsahu 80-200 slov v angličtině; nejedná se však o překlad
% zadání diplomové práce

Keywords:
Czech, Russian, Machine translation
% 3 až 5 klíčových slov v angličtině

\vss}

\newpage

%%% Strana s automaticky generovaným obsahem diplomové práce. U matematických
%%% prací je přípustné, aby seznam tabulek a zkratek, existují-li, byl umístěn
%%% na začátku práce, místo na jejím konci.

\openright
\pagestyle{plain}
\setcounter{page}{1}
\tableofcontents

%%% Jednotlivé kapitoly práce jsou pro přehlednost uloženy v samostatných souborech
\chapter*{Introduction}
\addcontentsline{toc}{chapter}{Introduction}
Slavic  family of languages consists of around 12 languages, de\-pen\-ding on the exact definition.\footnote{See for example \cite{sussex2011slavic} or \cite{siewierska1998overview} --  the final language count depends on whether Upper and Lower Sorbian are taken as different languages or not; whether Kashubian is a separate language, or just a variety of Polish; whether Serbian and Croatian are treated as separate languages or not}
They are usually divided into three sub-groups - South Slavic, West Slavic and East Slavic. While languages inside each of those groups are mutually intelligible, intelligibility accross those groups is lower; still, the Slavic languages in general have very similar morphology, syntax and vocabulary.

Czech language belongs to the West Slavic group, while Russian belongs to the East Slavic group.  Therefore, they are not fully mutually intelligible -- but, as noted, as Slavic languages, they are in some respects still very similar. 
This creates a great environment for building machine translations tools -- the languages are not too close to make the task meaningless, while we can still take an advantage of the closeness.

For political reasons, significant efforts were directed towards building an automatic translation system between Russian and Czech, mainly during mid-eighties with the RUSLAN system. The efforts were terminated in 1990 in the final phases of development.\footnote{See for example \cite{recycled}, \cite{hajic1987}, \cite{olivaruslan}} 

After the velvet revolution, political dependence on Russian-speaking countries were weakened, while the economical and cultural dependence on English-speaking countries were strengthened. Therefore, the efforts of automatic translation were more focused on English. 

In the process, general direction of machine translation efforts shifted from rule-based systems to more statistical-based systems. This allows us to compare older, rule-based approaches to the more modern statistical approaches.

In this thesis, I am trying to describe the rich history of automatic Czech-to-Russian translation, then experiment with the available tools -- both historical and more recent ones -- on a given set of data, and finally compare the results.

\chapter{Brief history of Czech and Russian}

\section{Proto-Slavic}

The Indo-Europeans occupied European lands from approximately 4000 BC. Their language, Proto-Indo-European (PIE), is not known and can only be linguistically reconstructed. 

In approximately 2000 BC, a Proto-Slavic (PS) language started to emerge. Until approximately 400 AD, the language was fairly uniform and the land occupied by Slavs was broadly coherent, stretched approximately from Oder river (Odra in Polish) to Dnieper (Днепр in Russian).\footnote{See \cite{sussex2011slavic}, pg.19 and further} 

In this part, I am describing some of the general changes from PIE to PS.
I am focusing mainly on syntactical changes -- not because there are no morphological, phonological or lexical changes, but because there are too many to cover here.

If a reader is interested in a further study of a history of Slavic languages and their difference in total, I can only recommend \cite{sussex2011slavic}. \footnote{Unless otherwise noted, information in this chapter are taken from this book.}

\subsection{Cases}
\jednatabulka{korintskym} { |l|X |X | X | }
{
         \hline
genitive &
invocant nomen \textbf{Domini} nostri Iesu Christi
&
   who call on the name \textbf{of} our \textbf{Lord} Jesus Christ         
&
všem, kdo na jakémkoli místě vzývají jméno našeho společného \textbf{Pána} Ježíše Krista
\\

    \hline
ablative
&
gratia vobis et pax a Deo Patre nostro et \textbf{Domino} Iesu Christo
&
Grace and peace to you \textbf{from} God our Father and \textbf{the Lord} Jesus Christ
&
Milost vám a pokoj od Boha, našeho Otce, a od \textbf{Pána} Ježíše Krista \\
    \hline
} {Demonstrating ablative and genitive conflation on 1 Corinthians} 


PIE had seven cases - nominative, accusative, genitive, dative, instrumental, locative, ablative and vocative. 
In PS, ablative and genitive were conflated into just genitive. 

To illustrate, I have added a simple comparison of old latin and Czech. 

Latin retained both ablative and genitive. Latin genitive of \emph{dominus} (lord, master) is \emph{dominī}, latin ablative of the same word is  \emph{dominō}. Both these forms were used in the very beginning of 1~Corinthians in Latina Vulgata (latin version of The Bible).

In the Table~\ref{tabulka:korintskym}, I have shown the translation to both Czech and English. We can see that both cases are inflated in Czech as genitive \uv{(našeho) pána}.\footnote{Bible sources: \cite{latinavulgata}, \cite{bibleniv}, \cite{bible21}. Note that English and Czech translations are not actually translations from Latin, but from more primary sources, but it will suffice for this simple comparison.}

\subsection{Numbers}
\jednatabulkan{kravy} { |l|l |l | l | }
{
         \hline
singular &
one cow
&
ena krava
&
jedna kráva
\\
   \hline
dual (in Slovenian) &
two cows
&
dve kravi
&
dvě krávy
\\
   \hline
plural &
three cows
&
tri krave
&
tři krávy
\\


    \hline
} {Demonstrating duals on Slovenian} 

PIE had three numbers, singular, plural and dual. In PS, dual slowly disappeared, while still retaining some of its usage.

However, dual disappeared in most of Slavic languages later (including Czech and Russian), leaving only traces in the grammar. One of the languages where dual remained is Slovenian.

To illustrate, I have added a comparison of \uv{one cow}, \uv{two cows} and \uv{three cows} in Slovenian and Czech in Table~\ref{tabulka:kravy}.

In Czech, dual was retained in declensions of several words, like \uv{hands}; in Russian, the dual \uv{рукама} survives in some dialects, but is generally incorrect.\footnote{See \cite{offord1996using}, page 18.}

\subsection{Genders}
PIE had three genders, masculine, feminine and neutral. PS retained these genders, adding categories Personal and Animate.

\subsection{Tenses}
PIE had six tenses: present, future, aorist, imperfect, perfect and pluperfect. 

The tenses were somehow retained in PS; however, future, perfect and pluperfect were in many cases re-formed analytically --- creating more complex forms with auxiliary verb and either an infinitive or past participle -- for example, \uv{budu zpívat} (I will sing) in Czech, or \uv{буду петь} in Russian.

\subsection{Moods}
PIE had four moods: indicative, subjunctive, optative and imperative. In PS, imperative forms were replaced by the optatives, and subjunctive mood slowly became conditional.

\subsection{Voices}
\jednatabulkan{voices} { |l|l |l | l | }
{
         \hline

active & νιπτω & I wash (someone) & myji (někoho) \\ \hline
medium & νίπτομαι & I wash (myself) & myji se \\ \hline
passive & νίπτομαι & I am washed (by somebody) & jsem myt \\ \hline


} {Voices in Classic Greek vs. Czech} 

PIE had an active and a mediopassive voice. 

In PS, this was refined as a reflexive and a non-reflexive voice, with the addition of a passive voice.

To illustrate, I have added an example of Classic Greek that still retained the mediopassive voice, in Table~\ref{tabulka:voices}.\footnote{See for example \cite{greek1}, \cite{greek2}}

\subsection{Aspects}
\jednatabulkan{nosim} { |l|l |l |  }
{
         \hline
imperfective determinate &
нести́
&
nést
\\
   \hline
imperfective indeterminate &
носи́ть
&
nosit
\\
   \hline
perfective& 
понести 
&
ponést

\\


    \hline
} {Aspects in Czech and Russian} 

PIE had distinction between two aspects -- eventive and stative; eventive aspect being further divided into perfective and imperfective aspect.\footnote{See \cite{ringe2008proto}, page 24}

In PS, the stative aspect is degramatized\footnote{See \cite{andersen2013origin}}; however, the perfective/imperfective distinction became more important in PS than in other Indo-European languages. Also, the imperfective motion verbs were further split into determinate and non-determinate.

The determinate/non-determinate and perfective/imperfective distinction is retained in both Czech and Russian. See the Table~\ref{tabulka:nosim} and note, how hard would be to correctly translate the distinction into English.


\section{Drifting of Czech and Russian}
In about 5th century AD, the relative unity of Slavic started to break up, caused mainly by migration, which was partly caused by expansion to the north and east by the Eastern Slavs, partly by dissolution of Roman and Hun empires and the resulting vacuum in Central Europe.

One group of Slavs moved westwards, reaching approximately what is now Poland, Czech Republic and parts of Germany, creating the so-called West Slavic language group, which Czech language is a member of. The other group moved south to Balkan, creating a South Slavic language group, which I won't discuss further in this thesis. Eastern Slavs moved north and east, eventually creating an Eastern Slavic language group. Slavic was divided into those three groups at about 10th century AD.

The differences between the languages themselves are actually less syntactical and are more phonological, morphonological and morphological, with some broad-scale lexical changes. For that reason, I am not describing the linguistic differences of the languages any further, since it would be, again, not in a scope for this thesis. Again, I can recommend the book \cite{sussex2011slavic} for anybody interested in the deeper differences.



\chapter{Experiment setting}
\section{Statistical vs. rule-based translation systems}

Machine translation (MT) systems has historically used many different approaches. One way of classifying the approaches is on the axis of rule-based vs. statistical.

In general, we can re-use the descriptions, used by \cite{bojar}, which is as follows:
\begin{itemize}
\item rule based MT systems:
\begin{itemize}
\item use analysis, transfer and synthesis steps
\item rule-based machine translation systems use formal grammars
\item rule-based machine translation systems use hand-made dictionaries
\item has linguistic information hard-coded and therefore aren't language-agnostic
\end{itemize}
\item statistical MT systems
\begin{itemize}
    \item use more variants of outputs, rank them with some score, and choose the best one
    \item train internal dictionaries from big parallel data
    \item the core of the system is compact and its inner working less obvious
    \item uses statistics instead of linguistic rules and therefore are language-agnostic
\end{itemize}
\end{itemize}

However, with some of the modern system, the distinction is not as clear-cut as we would like, for the purpose of our comparison. 
For example, statistical MT systems like Moses can get significantly better with some added linguistic information; 
on the other hand, systems like TectoMT, which can sometimes be classified as more rule-based, actually have big modules more or less based on statistics.

%In this thesis, I will present several existing solutions to Czech-Russian machine translation, describe their position on the aforementioned axis, present their history and discuss their results on a testing data.
\section{Development and testing data}

To compare several machine translation systems, we need some consistent data to show the deficiencies and/or improvements of different MT systems.

In general, we want to produce a system that's general enough and performs well on unseen sentences; for that reason, we should have -- in the case of statistical MT systems -- separated test set and training/development set.\footnote{As noted in for example \cite{koehn2010statistical}, p. 274, or \cite{bishop}, pg. 6}

We have tried to do that in our experiments with our data models. However, some systems evaluated in this thesis are only black-box systems (Google Translate, PC Translator). Especially with Google Translate, we cannot rule out the possibility that they already use some parts of the test/development data.






\chapter{Black-box systems}

\section{PC Translator}

\subsection{Description}

PC Translator is a commercial translation system from a Czech company LangSoft (\url{http://www.langsoft.cz/translator.htm}). PC Translator can translate several language pairs with Czech on either source or target side.

Authors of PC Translator don't publish any papers or other literature about the system -- what can we tell about its functionality is gathered only from its promotional website and from the experiments with the software itself.

PC translator seems to be purely rule-based. The system seems to work in following steps:

\begin{itemize}
\item some (probably rule-based) morphological analysis of the source language
\item translation of the lemmas from source language to target language by searching in a large dictionary
\item some synthesis of morphological information and the translated lemma
\end{itemize}

The system doesn't seem to do any kind of reordering. It also doesn't seem to do any analysis on a deeper level, like sentence constituents. Some of the phrases in the dictionary are longer than one word, but not too many of them.

One of the advantages of PC Translator is its large dictionary -- however, the dictionary is sometimes choosing very odd and inprobable choices when disambiguating between more possible translations. For example, the English sentence \uv{I like dogs} is translated as \uv{Mám rád kleště}, because the term \uv{dog} can be also translated as \uv{kleště}\footnote{from Collins' Dictionary: \uv{dog -- 5. a mechanical device for gripping or holding, esp one of the axial slots by which gear wheels or shafts are engaged to transmit torque}}.

According to its marketing materials, PC Translator v14 uses a Czech-Russian dictionary with above 650.000 words.

\subsection{Experiments}
We found out it's not easy to automatize translating with PC Translator. Its GUI is suited for translating by hand, sentence-by-sentence, but not for automated translation of thousands of sentences. Also, by definition, Windows GUI is harder to automate on Linux machine from a script.

However, we were able to do some work-around, with the help of VMWare Player virtualization software (\url{http://www.vmware.com/cz/products/player}) and AutoHotkey GUI scripting software (\url{http://www.autohotkey.com/}). Our workflow therefore is:

\begin{itemize}
\item on Linux machine, encode the source from UTF-8 to windows-friendly encoding
\item encode the source as HTML code
\item start a virtual machine with PC Translator pre-installed
\item on the start of the virtual machine, run AutoHotkey script from an outer-machine folder (thanks to VMWare shared folders and Windows Startup scripts)
\item via this AutoHotkey script, run PC Translator and click on \uv{translate file} feature 
\item translate the HTML file (also shared in the VMWare shared folder)
\item turn off the virtual machine
\item turn the file back from HTML and Windows encodings back to UTF-8
\end{itemize}

The HTML part is needed because PC Translator had some problems with translating ordinary text files, plus we can pair the translated sentences better thanks to \texttt{id} parameters in \texttt{div} tags.

We used the newest version of PC Translator available at the time, which is PC Translator v14.



\subsection{Results}

\jednatabulkan{bleutrans} { |r|r|r | }
{
\hline
&
intercorp&
WMT\\ \hline
BLEU & 5.02\% $\pm$ 0.01\%
&
6.74\% $\pm$ 0.32\%

\\ \hline
}{PC translator BLEU}

\jednatabulka{transsent1} { |X|X|X | }
{
\hline
Je mi teprve čtyřiadvacet let a nemohu prožit celý svůj život s legitimací invalidy práce a potloukat se po nemocnicích , když vím , že je to marné .   &   Есть мне только двадцать четыре рейс и не могу пережив весь свой жизнь с удостоверение личности инвалидами труд и болтаться по больницах , когда знаю , что это бесполезные .   &   Est` mne tol`ko dvadczat` chety're rejs i ne mogu perezhiv ves` svoj zhizn` s udostoverenie lichnosti invalidami trud i boltat`sya po bol`niczax , kogda znayu , chto e`to bespolezny'e .\\ \hline
Generál chodil po pokoji sem a tam , kouře svou pěnovku .   &   Генерал ходил по комнате туда - сюда , дыма свою пенковая трубка .   &   General xodil po komnate tuda - syuda , dy'ma svoyu penkovaya trubka .\\ \hline
" Možná že žádné brilianty neexistují ? "   &   " возможно никакие бриллианты отсутствовать ?    &   " vozmozhno nikakie brillianty' otsutstvovat` ? \\ \hline
Nesměl při tom udělat chybu , vyžadovalo to stejnou přesnost , jako když se zaměřuje dělo .   &   Обмен причём ошибиться , требовало то такой же аккуратность , как когда специализируется пушка .   &   Obmen prichyom oshibit`sya , trebovalo to takoj zhe akkuratnost` , kak kogda specializiruetsya pushka .\\ \hline
S hrdostí vzpomněl , jak snadno dobyl kdysi srdce krásné Heleny Baurové .   &   С гордостью вспомнил , как ходко захватил когда - то сердце красивое Елена Баурове .   &   S gordost`yu vspomnil , kak xodko zaxvatil kogda - to serdce krasivoe Elena Baurove .\\ \hline

}{PC Translator demonstration - InterCorp}
\jednatabulka{transsent2} { |X|X|X | }
{
\hline
Thiago Silva, který patří k nejlepším obráncům na světě, taky umožňuje ostatním vedle sebe růst.   &   Тhиаго Силва, какой принадлежать к лучшим защитник в мире, тоже даёт возможность остальным плечом к плечу рост.   &   Thiago Silva, kakoj prinadlezhat` k luchshim zashhitnik v mire, tozhe dayot vozmozhnost` ostal`ny'm plechom k plechu rost.\\ \hline
"Dávali mi pět let života a už je to sedm," říká bez emocí na svém lůžku v domě pro paliativní péči Victor-Gadbois v Beloeil, kam přijel předešlý den.   &   "давали мне пять лет жизни и уже это семь," говорит минус эмоций в своем гнезде в доме для башка опеку Victor-Гадбоис зажечься Бєлоєил, куда доехал предшествующий день.   &   "davali mne pyat` let zhizni i uzhe e`to sem`," govorit minus e`mocij v svoem gnezde v dome dlya bashka opeku Victor-Gadbois zazhech`sya Bєloєil, kuda doexal predshestvuyushhij den`.\\ \hline
Opatrnost je ovšem na místě například na některých přemostěních, kde může být povrch namrzlý a kluzký.   &   Осторожность есть конечно на месте например на некоторых перекрытие, где может быть покрытие замёрзший и сальный.   &   Ostorozhnost` est` konechno na meste naprimer na nekotory'x perekry'tie, gde mozhet by't` pokry'tie zamyorzshij i sal`ny'j.\\ \hline
Prostě je ignoruji.   &   Запросто есть игнорирую.   &   Zaprosto est` ignoriruyu.\\ \hline
Podle doktorky Christiane Martelové není quebecký zdravotnický systém dostatečně výkonný, aby zajistil přístup všech osob ke kvalitní paliativní péči, než bude možno souhlasit s provedením eutanazie.   &   По врачи Christiane Мартєлове нет qуєбєцкэ медицинский комплекс достаточно производительный, чтобы обеспечил подход всех» личностей к качественный башка опеку, нежели можно будет согласиться с проведением єутаназиє.   &   Po vrachi Christiane Martєlove net quєbєczke` medicinskij kompleks dostatochno proizvoditel`ny'j, chtoby' obespechil podxod vsex» lichnostej k kachestvenny'j bashka opeku, nezheli mozhno budet soglasit`sya s provedeniem єutanaziє.\\ \hline


}{PC Translator demonstration - WMT2013}

We can see the results are not ideal. The sentences are evidently translated word-by-word, without any regard for the whole sentence, or the correct word order.

Because the languages are similar, the morphological information is sometimes transferred correctly -- for example, \uv{s hrdostí vzpomněl} -> \uv{с гордостью вспомнил} -- but sometimes, very incorrectly -- for example, genitive \uv{překrásné} is not with the correct ending as \uv{-ой}, but incorrectly as \uv{-ое}.

It's obvious the dictionary is quite big and complete, so a lot of the words were translated somehow correctly; however, some of the words had strange ?????????


\section{Google Translate}
\subsection{Description}
Google Translate is a popular free online translation service by Google, an American web search giant (\url{http://translate.google.com}). 
Although Google is producing many academic papers on machine translation, the whole system is still proprietary and we cannot fully inspect it, as in the case of PC Translator.

Google Translate uses mostly statistical approach to machine translation, see for example \cite{och}\footnote{F. J. Och is a head of Google Translate group in Google}. Its results are often seen as \uv{state-of-the-art} in machine translation.

However, thanks to its purely statistical approach, it either needs huge amounts of data for every language pair, or it needs to use so-called \uv{pivot languages}\footnote{See for example \cite{koehn2010statistical}} -- in the case of Google Translate, it's usually English; specific English word order and English idioms are then re-translated into the target language and sometimes introduce downright wrong translations.

\subsection{Experiments}
To automate Google Translate, we cannot use the website itself, simply because pasting tens of thousands of lines into a browser window usually crashes the browser itself.

There are some workarounds around this, such as \uv{faking} browser environment using some automation tools and/or libraries, but we used more stable option.

Google Translate, apart from being a website, has a paid translation API\footnote{\url{https://developers.google.com/translate/?hl=cs}}. We figured out it's not too expensive for our testing purposes, so we ended up paying for the API.\footnote{The cost is measured per character on the source side. We used about 3 million characters and paid about 60 dollars. This is rather high for any repeated experiments, but not that high for one-time translation.}

We used an unofficial Java library for Google Translate API, called prosaically \uv{google-api-translate-java}\footnote{https://code.google.com/p/google-api-translate-java/}.

The tests were done on 3rd May, 2014.\footnote{I think it's important to note the date of the tests, because the quality of online services might change overtime.}

\subsection{Results}


\jednatabulkan{bleugoog} { |r|r|r | }
{
\hline
&
intercorp&
WMT\\ \hline
BLEU & 8.79\% $\pm$ 0.15\%
&
17.96\% $\pm$ 0.32\%

\\ \hline
}{Google translate BLEU}

The first thing to note is the radically different results for Intercorp and WMT corpus. This disprepancy is not repeated in our other experiments, where both sets have BLEU that's closer to each other.

This might be caused by two reasons. One reason might be that the complete WMT data are simply already included in Google's training data, while Intercorp is not. It's not that far-fetched theory: WMT is one year old already, and Google is one of WMT sponsors. On the other hand, in this case, the results should be even better than they actually are -- Google's results are, after all, worse, than our Moses experiments in chapter ??.

Other reason might be that Google is domain specificity of Google's training data. It's possible that Google is generally using far more newspaper data for training than just fiction, which might disadvantage Google's system in Intercorp case. Since we don't know Google's training data structure, we can't know for sure.



\jednatabulkam{googsent1} { |X|X|X | }
{
\hline

Je mi teprve čtyřiadvacet let a nemohu prožit celý svůj život s legitimací invalidy práce a potloukat se po nemocnicích , když vím , že je to marné .   &   Я всего двадцать четыре года, и я живу свою жизнь с купоном недействительным работы и болтаться в больнице, зная, что это бесполезно.   &   Ya vsego dvadczat` chety're goda, i ya zhivu svoyu zhizn` s kuponom nedejstvitel`ny'm raboty' i boltat`sya v bol`nice, znaya, chto e`to bespolezno.\\ \hline
Generál chodil po pokoji sem a tam , kouře svou pěnovku .   &   Генеральный ходил по комнате взад и вперед, его дым пены.   &   General`ny'j xodil po komnate vzad i vpered, ego dy'm peny'.\\ \hline
" Možná že žádné brilianty neexistují ? "   &   "Возможно, не алмазы там?"   &   "Vozmozhno, ne almazy' tam?"\\ \hline
Nesměl při tom udělat chybu , vyžadovalo to stejnou přesnost , jako když se zaměřuje dělo .   &   Он не мог ошибиться в этом, это требуется такой же точности, как при фокусировке пушки.   &   On ne mog oshibit`sya v e`tom, e`to trebuetsya takoj zhe tochnosti, kak pri fokusirovke pushki.\\ \hline
S hrdostí vzpomněl , jak snadno dobyl kdysi srdce krásné Heleny Baurové .   &   Мы с гордостью вспомнил, как легко покорил когда-то сердце прекрасной Елены Баур.   &   My' s gordost`yu vspomnil, kak legko pokoril kogda-to serdce prekrasnoj Eleny' Baur.\\ \hline


}{Google Translate demonstration -- Intercorp}


\jednatabulkam{googsent2} { |X|X|X | }
{
\hline

Thiago Silva, který patří k nejlepším obráncům na světě, taky umožňuje ostatním vedle sebe růst.   &   Тьяго Силва, один из лучших защитников в мире, также позволяет другой параллельный рост.   &   T`yago Silva, odin iz luchshix zashhitnikov v mire, takzhe pozvolyaet drugoj parallel`ny'j rost.\\ \hline
"Dávali mi pět let života a už je to sedm," říká bez emocí na svém lůžku v domě pro paliativní péči Victor-Gadbois v Beloeil, kam přijel předešlý den.   &   "Они дали мне пять лет жизни, и это семь", говорит он без эмоций на его постели у себя дома для паллиативной помощи Виктор-Гадбуа в Beloeil, куда они прибыли в предыдущий день.   &   "Oni dali mne pyat` let zhizni, i e`to sem`", govorit on bez e`mocij na ego posteli u sebya doma dlya palliativnoj pomoshhi Viktor-Gadbua v Beloeil, kuda oni priby'li v predy'dushhij den`.\\ \hline
Opatrnost je ovšem na místě například na některých přemostěních, kde může být povrch namrzlý a kluzký.   &   Внимание, однако, находится в месте, например, некоторые перемычки, где поверхность может быть ледяной и скользкий.   &   Vnimanie, odnako, naxoditsya v meste, naprimer, nekotory'e peremy'chki, gde poverxnost` mozhet by't` ledyanoj i skol`zkij.\\ \hline
Prostě je ignoruji.   &   Просто игнорируйте их.   &   Prosto ignorirujte ix.\\ \hline
Podle doktorky Christiane Martelové není quebecký zdravotnický systém dostatečně výkonný, aby zajistil přístup všech osob ke kvalitní paliativní péči, než bude možno souhlasit s provedením eutanazie.   &   По словам доктора Кристиана Martel Квебеке система здравоохранения не является достаточно мощным, чтобы обеспечить доступ для всех людей на высококачественной паллиативной помощи, прежде чем он может согласиться проводить эвтаназию.   &   Po slovam doktora Kristiana Martel Kvebeke sistema zdravooxraneniya ne yavlyaetsya dostatochno moshhny'm, chtoby' obespechit` dostup dlya vsex lyudej na vy'sokokachestvennoj palliativnoj pomoshhi, prezhde chem on mozhet soglasit`sya provodit` e`vtanaziyu.\\ \hline


}{Google Translate demonstration -- WMT2013}

NĚJAKÉ HODNOCENÍ VÝSLEDKŮ - já osobně to nedokážu vůbec zhodnotit.

\section{Microsoft Bing Translator}
\subsection{Description}
Another online service that we decided to try is Microsoft Translator/Bing Translator. (In Microsoft's own materials, the system is usually called Bing Translator when referring to the website and Microsoft Translator when referring to the API.)

Microsoft Translator is very similar to Google Translate -- it is an online website with an easy GUI, and an additional paid API. Again, the team occasionally publishes some scientific papers, but is again otherwise proprietary.

In different experiment, we found out (non-scientifically), that for some language pairs, Microsoft Translator does more post-editation, that seemed a little rule-based (for example, better verb separation in English-to-German translation). Therefore, we decided to try it alongside Google Translate in our experiments.


\subsection{Experiments}
Again, we used Microsoft's Translator API (marketed as part of Windows Azure).

The API is slightly more complicated than Google's API because of the auto-expiring token, but we used the example PHP script from the API documentation\footnote{\url{http://msdn.microsoft.com/en-us/library/hh454950.aspx}}.

The pricing is slightly different in Microsoft Translator than in Google Translate, but in general is slightly cheaper. First 2 million letters are for free, next 2 million are for about 40 US dollars.

\chapter{Moses}

\section{General description}
Moses is an open-source machine translation toolkit with GPL licence, developed as a successor to the closed-source Pharaoh system.\footnote{See for example \cite{mosespaper} or \cite{moseslink} -- but Moses is used so often and so extensively that many other papers could be found}

The system is very modular and very customizable, which makes it a bit harder to describe. In this section, I will try to describe our Moses set-up; bear in mind, however, that a completely different set-up is also possible.


\subsection{Pipeline overview}
At the start, we have a bilingual corpora of a given language pair, and bigger monolingual corpora of the target language.

The bilingual corpora has to be prepared by aligning the sentences, so every sentence has exactly one translation. We describe our corpora in the part ???

The sentences are then word-aligned, which means pairing words to their translations. We are using MGIZA++\footnote{See \cite{mgiza}}.

From this word alignment, Moses learns a so-called \emph{phrase-based translation model}.

From the monolingual corpora, we then learn a statistical language model. We use SRILM language model\footnote{See \cite{srilm}}.

Moses is then used for so-called \emph{decoding} of the information from both the language model and the translation model, which choses the best possible translation, using things like beam-search.

However, for the best translation, we need to tune Moses parameters for optimal results. This is done using so-called \emph{minimum rate error rating} -- or MERT for short, which is tuning the parameters on a small separate development set.

After MERT tuning, we finally have working language model, translation model and Moses parameters, which is our complete translation system.

To reiterate, this is our Moses pipeline
\begin{pitemize}
\item getting sentence-aligned parallel Czech-Russian corpus, plus Russian monolingual corporus
\item world-alignment on parallel corpus
\item creating phrase-based translation model
\item creating Russian language model
\item tuning the parameters for Moses decoder
\end{pitemize}

For managing input and output from the various steps, we use \emph{eman} system, which we transformed a little.


\section{Our Moses setup}
In this section, I am describing our exact setup and our data.

\subsection{Parallel data}
\subsubsection{Intercorp}
In general, we described Intercorp corpus in \ref{intercorp_p1}.


We created the Moses system before we started evaluating the system on the evaluation corpus, as described in Section \ref{intercorp_p1}. 
The data we had access to were more limited at that time: they were from the previous version of Intercorp, didn't have all the metadata, such as the publishing dates and information about the original translation source, and were shuffled. However, we can reconstruct the metadata by comparing this shuffled data with the newest version of intercorp.


TODO: popis intercorp dat

\subsubsection{Subtitles}
Another set of data that we used were subtitles from OpenSubtitles database.

In a separate FilmTit project\footnote{???citovat ročníkový projekt/dokumentaci k němu?}, we tried to make a machine translation project for subtitle translation from English to Czech.

For that project, we were given access to the set of subtitles from the server OpenSubtitles (\url{http://opensubtitles.org}). This dataset was not a pair of aligned sentences, not even a pair of aligned \emph{files}; we were given just a set of SRT files, and a table which paired each of those files with a movie (identified by IMDB number). Each movie usually has more subtitle files.

Subtitle files have the sentences paired with timestamps.

From a set of SRT files paired with a movie, we selected just one Czech and just one English SRT file which we find most similar, based on the timestamps.
From the pair of the files, we then extract the sentences that have the most similar timestamps.

These two pairings -- pairings of subtitle files and pairing of the actual lines -- are non-trivial task, and require a \uv{tolerance} -- how different can the times of a sentence be to be still paired together.

Higher tolerance produces bigger corpus with more errors, while lower tolerance produces smaller, but more correct corpus.

When we were experimenting on the afforementioned project, we found out, that the best results (tested both on another movie subtitles and a different corpus) are -- without exception -- with \emph{bigger corpus} and \emph{higher tolerance}. Even when that introduced a lot of incorrect sentence pairs, the overall results were still better with the biggest possible corpus.

For Czech-to-Russian experiments, we were given another, similar set of subtitels from OpenSubtitles. The set was naturally much smaller than with English and Czech, but we were still able to use the same algorithms to build a parallel corpora.

We again used the highest possible tolerance, and therefore surely introduced a lot of errors. Unfortunately, for a lack of time, we weren't able to replicate the experiments for the ideal tolerance here, however, we hoped that the results would be similar than in English-to-Czech translation, that is, the biggest corpus resulting in the best translation.

\subsubsection{UMC corpus}
UMC (UFAL Multilingual Corpus) is a Czech-English-Russian corpus\footnote{See \cite{umc}}. The data are downloaded from Project Syndicate, a Prague-based non-profit news organization, translating news and opinions from around the world.

TODO: počty vět a tokenů, jediná věc, co o UMC můžu napsat

\subsubsection{Wiki titles}
We also extracted all of the titles from the Czech and Russian wikipedia, that correspond to each other.

Those are usually only noun phrases and the main word is usually in nominative singular, so the morphology isn't that rich -- however, we hope that the model will learn some phrases needed for the translation of the named entities.

\subsection{Monolignual data}
\subsubsection{News Crawl}
The largest part of our monolingual data is corpus from WMT workshops that they call \uv{News Crawl}.

According to \cite{wmt_findings_2009}, WMT workshop has been continously crawling web articles since 2007 for making test sets. This allowed them to make a big, randomized corpus from all these sources.

The corpus is categorized by year, and we treat each year as its own corpus for the interpolation.
\subsubsection{Common Crawl}
Common Crawl is publicly available web crawl\footnote{From Wikipedia -- \uv{A Web crawler is an Internet bot that systematically browses the World Wide Web, typically for the purpose of Web indexing. }; web crawl is then a result of such a web crawler} -- \url{http://commoncrawl.org/}.

As described in \cite{commoncrawl}, group of researchers tried to extract parallel data from this web crawl. One of the language pairs was English-Russian and the result is publicly available on WMT site. We used the Russian side of the corpus.

However, the quality of this corpus is very discutable. Because it contains data downloaded from the \uv{raw web}, it often has sentences in different languages, sentences in machine-translated Russian, random UTF-8 symbols, random HTML data, some code, and so on.

This was one of the reasons why we decided to use linear interpolation as discussed above -- hoping, that the tuning algorithm will automatically \uv{find the right balance} between the language models.

\subsubsection{Yandex}
%Yandex is a Russian portal and search service, that can be described like a Russian mix between Google and Yahoo!. One of its services is a free translation service
Yandex was already described in the chapter ??. Apart from providing free translation API, Yandex also provides an English-Russian parallel corpus.

We used the Russian part of this corpus for language modelling.

\subsection{Managing experiments}
For managing the steps described further, such as training models, we need some overaching system -- steps variously fail, don't compile, don't fit in memory, etc., plus we would like to reuse partial results in more experiments.

Moses itself has built-in perl-based experiment managment system, called prosaically Experiment Management System (EMS). However, this system is not very widely used in on UFAL.

Instead of EMS, we use instead another perl-based tool called eman (experiment manager). Eman is described well in \cite{eman} or at its website, \url{http://ufal.mff.cuni.cz/eman}.  

Eman breaks down experiment into so-called \uv{steps}. Step encapsulates an atomic part of an experiment and can be in one of a few various states. More importantly, step can be dependent on various other states; if a step fails, all steps dependent on it automatically fail. The whole experiment is then just another step, dependent on all the necessary substeps.

Step is represented by a directory in a playground directory. Step is created by copying a script, called \uv{seed}, from a library of seeds, to a new directory.

In our opinion, eman itself is well written, while we found the seeds themselves hard to read, too repetitive, and with large amount of code copied and pasted over. We tried to rewrite the seeds as perl modules, so the reusability is higher.

I am not personally sure if this effort was successful. We decided to use MooseX::Declare, which seemed to us like a modern way to write modules in perl. 

Unfortunately, MooseX::Declare is using very difficult-to-understand perl concepts and source code transformations, and as a result, it takes long to run and, perhaps worse, returns very confusing and undecypherable errors. So as a result of our refractoring, we have seeds that have code that's probably easier to read and refractor, but on the other hand, it's slow and produces very opaque errors.

\subsection{Word alignment}
For word alignment, we are using MGIZA++\footnote{See \cite{mgiza}}, which is a GPL toolkit based on GIZA++\footnote{See \cite{giza}}, which is itself based on models, sometimes called IBM Model 1 to IBM Model 5\footnote{See \cite{ibm}}, which is itself based on expectation–maximization algorithm (EM).

IBM Models and the underlying EM algorithms are explained perfectly in Chapter~4 of \cite{koehn2010statistical} or in those slides by the same author -- \url{http://www.inf.ed.ac.uk/teaching/courses/mt/lectures/ibm-model1.pdf}.

GIZA++ is an implementation of those models. MGIZA++ is just its multi-threaded variant, which makes the word alignment slightly faster.

\subsection{Phrase-extraction}
In this step, Moses takes the word alignment from the previous step and learns a so-called \uv{phrase table}.
Unlike word alignment, phrase extraction spans multiple words on every side in so-called \uv{phrases}.

Phrase table consists of list of phrases, their probabilities in both ways of translation, and their lexical weighting -- lexical weighting is the probability of the translated phrase counted by individual word pairs. The exact meaning of the numbers is well explained in \cite{koehn2003}.

The phrase-table - defines a so-called \uv{translation model}. 

\subsection{Language model}
Language model is a part of the system, that tries to model the probability of a target language sentence alone. It's trained on a monolingual corpus.

We use SRILM, which is an open source language modeling toolkit. (Although it's open-source, it uses its own license, that allows free use only for non-commercial and educational purposes.) Current status of SRILM is described in \cite{srilm}, original design is described in \cite{srilm_old}. 

SRILM uses several models, one of them is n-gram word model, described well in \cite{koehn2010statistical}\footnote{chapter 7}. We use n-grams model to the order 3 with words and order 5 with tags (see ???). We smooth the models with Kreser-Ney smoothing with Chen and Goodman's modification\footnote{See \cite{chen} and \url{http://www.speech.sri.com/projects/srilm/manpages/ngram-discount.7.html}}. 


\subsection{Language model interpolation}

We had more than one monolingual Russian corpora, but we weren't sure of how high quality each of them were and how helpful it would be. For this reason, we used so-call interpolation (also called mixing).

Linear interpolation in general is described for example in \cite{gutkin}. On a separate heldout data, set of \emph{lambdas} are trained -- the resulting probabilities are then just the individual probabilities, multiplied by the lambdas and summed.

Linear interpolation is supported by undocumented script in Moses \texttt{interpolate-lm.perl}, which in turn uses SRILM's undocumented AWK script \texttt{compute-best-mix.gawk} and SRILM's \texttt{ngram} with \texttt{-mix-lm} option\footnote{See \url{http://www.speech.sri.com/projects/srilm/manpages/ngram.1.html}}.

We used a linear interpolation instead of log-linear interpolation simply because we didn't notice the option until later in the project.

\subsection{Factored translation}
\subsubsection{Overview}

The pipeline, described in the previous section, translates phrases from the source language directly to the target language.

However, with morphologically rich languages such as Russian or Czech, this can result in too low results. With so-called factored translation, we can add some morphological information while still keeping the main ideas of phrase-based translation.

With factored translation, we add some morphological (or other) information to either source, target or both -- for example, lemma or morphological tag -- this is called a \emph{factor}. Then, instead of making language models and/or translation models on the words alone, we train them on some combination of these factors and then somehow put them together.

\subsubsection{Factor experiments}
In a separate set of experiments, that predate the whole experiment described here, only on the UMC data described in the section ??, we compared several set-ups for factored translation.

We used a set-up described for example in \cite{backoff}. However, since we do not have Russian morhpology fully working, we used only the system described as \texttt{lemma backoff} -- with the exception of not translating to lemma. We were not using interpolated backoff, simply because regular backoff is easier to use with Moses.

The main model translates from a word form on the source side to word form and tag on the target side. The backoff model translates from a lemma (or a stem -- see below) to form and tag on the target side.

For Russian, we used TreeTagger software\footnote{\url{http://www.cis.uni-muenchen.de/~schmid/tools/TreeTagger/}, also see \cite{treetagger1} and \cite{treetagger2}} with a Russian parameter file\footnote{trained on a corpus created by Serge Sharoff, see \url{http://corpus.leeds.ac.uk/mocky/}}. TreeTagger is a closed-source software with a restrictive license, but for free for research purposes.

For Czech, we used tokenizer from UFAL project Treex (described further in section XXX) and morphological analyzer Morče (\url{http://ufal.mff.cuni.cz/morce/references.php}); however, we didn't actually use its output, as described further.

When we experimented further, we found out that using not lemma on the source side, but a very crude lexical stem -- just using the first 4 letters of a word -- gets better results\footnote{Using stems instead of lemmas is suggested for example in XX. However, their stems are more linguistically motivated, while we just crudely take first few letter}. When we tried the experiment with different lengths of the stem, we got the results, shown in the graph XX -- the lexical stem of the lengths XX was the best.

Therefore, we continued to use that in our further experiments, as for example in the WMT submissions\footnote{citace tu}.




???? výsledky, ale stejně to budu přehazovat











% Ukázka použití některých konstrukcí LateXu (odkomentujte, chcete-li)
% \include{example}


\chapter*{Conclusion}
\addcontentsline{toc}{chapter}{Conclusion}
TODO


%%% Seznam použité literatury
%%% Seznam použité literatury je zpracován podle platných standardů. Povinnou citační
%%% normou pro diplomovou práci je ISO 690. Jména časopisů lze uvádět zkráceně, ale jen
%%% v kodifikované podobě. Všechny použité zdroje a prameny musí být řádně citovány.

\def\bibname{Bibliography}
%\begin{thebibliography}{99}
\addcontentsline{toc}{chapter}{\bibname}


\printbibliography



%\bibitem{lamport94}
%  {\sc Lamport,} Leslie.
%  \emph{\LaTeX: A Document Preparation System}.
%  2. vydání.
%  Massachusetts: Addison Wesley, 1994.
%  ISBN 0-201-52983-1.

%\end{thebibliography}


%%% Tabulky v diplomové práci, existují-li.
\chapwithtoc{List of Tables}

%%% Použité zkratky v diplomové práci, existují-li, včetně jejich vysvětlení.
\chapwithtoc{List of Abbreviations}

%%% Přílohy k diplomové práci, existují-li (různé dodatky jako výpisy programů,
%%% diagramy apod.). Každá příloha musí být alespoň jednou odkazována z vlastního
%%% textu práce. Přílohy se číslují.
\chapwithtoc{Attachments}

\openright
\end{document}
